% Exact quarantined release candidate for P06.2.
% This file is not a live-book input before independent and production
% verification.  No novelty or priority claim is made.

\begin{proof}[Proof of \cref{prop:p062-rational-centred-traces}]
Fix \(a\in D^-_{\mathbf R}(K)\) and a polyradius
\(\mathbf r\) with \(0<r_i<R_i\).  Choose numbers
\begin{equation}\label{eq:p062-rho-above-r}
                         r_i<\rho_i<R_i
\end{equation}
for every \(i\), and put
\[
 \Omega_{\boldsymbol\rho}
   =\prod_{i=1}^n(-\log\rho_i,\infty),
 \qquad
 A_{a,\boldsymbol\rho}
   =\{0<|T_i-a_i|<\rho_i\text{ for every }i\}.
\]
The ultrametric inequality places \(A_{a,\boldsymbol\rho}\) inside
\(D^-_{\mathbf R}\).  Translation by \(a\) identifies it with the split
torus domain
\(\operatorname{trop}^{-1}(\Omega_{\boldsymbol\rho})\), whose
tropicalization and weight-one Gauss section are
\[
 F_a=(-\log|T_1-a_1|,\ldots,-\log|T_n-a_n|),
 \qquad
 \iota_a(u)=\eta_{a,(\mathrm e^{-u_1},\ldots,\mathrm e^{-u_n})}.
\]

Let \(\beta\) be a compactly supported smooth Lagerberg form of
bidegree \((n-1,n-1)\) on \(\Omega_{\boldsymbol\rho}\).  Tropicalization
of a split torus is proper.  Thus \(F_a^*\beta\) is a compactly supported
smooth test form on \(A_{a,\boldsymbol\rho}\).  Restricting the current
equation to this ordinary open analytic subspace, or equivalently
extending this test form by zero, gives
\begin{equation}\label{eq:p062-arbitrary-field-test-identity}
 \begin{aligned}
 0
 &=\bigl\langle\mathrm d'\mathrm d''[h],F_a^*\beta\bigr\rangle\\
 &=\int_{A_{a,\boldsymbol\rho}}
             h\,F_a^*(\mathrm d'\mathrm d''\beta).
 \end{aligned}
\end{equation}

Set \(\alpha=\mathrm d'\mathrm d''\beta\) and let
\(\mu_\alpha=(F_a^*\alpha)_{\mathrm{Rad}}\) be its signed Radon
measure.  Choose a compact \((\mathbf Z,-\log|K^\times|)\)-polytope
\(\Delta\Subset\Omega_{\boldsymbol\rho}\) containing
\(\operatorname{supp}(\alpha)\).  Such a polytope exists even when the
valuation is discrete: the divisible hull of the nonzero additive group
\(-\log|K^\times|\) is dense in \(\mathbf R\), and rational integral
inequalities therefore provide a compact box strictly inside
\(\Omega_{\boldsymbol\rho}\) around the given support.  The top-form
support theorem on a torus
\cite[Proposition~2.3.6]{BGJKK25} and the weight-one calibrated
integration formula
\cite[Proposition~11.1 and Example~11.7]{GJR26} give
\begin{equation}\label{eq:p062-arbitrary-field-support-push}
 \operatorname{supp}(\mu_\alpha)
       \subseteq\iota_a(\operatorname{supp}\alpha),
 \qquad
 (F_a)_*\mu_\alpha=\alpha_{\mathrm{Rad}}.
\end{equation}
These statements do not require \(K\) to be algebraically closed.  The
calibrated formula first gives the second equality against smooth
compactly supported multipliers.  Such multipliers are uniformly dense
on the common compact support, so uniqueness of signed Radon measures
gives the displayed equality.  Since \(F_a\circ\iota_a\) is the identity,
the support and pushforward identities imply
\begin{equation}\label{eq:p062-arbitrary-field-measure-section}
                      \mu_\alpha
                         =(\iota_a)_*\alpha_{\mathrm{Rad}}.
\end{equation}
The multiplier \(h\) is continuous, so
\eqref{eq:p062-arbitrary-field-test-identity} and
\eqref{eq:p062-arbitrary-field-measure-section} yield
\begin{equation}\label{eq:p062-arbitrary-field-trace-hessian}
  \int_{\Omega_{\boldsymbol\rho}}
       h(\iota_a(u))\,\mathrm d'\mathrm d''\beta(u)=0
\end{equation}
for every such \(\beta\).  Thus the continuous function
\[
                    \tau_a(u)=h(\iota_a(u))
\]
has zero distributional Lagerberg Hessian.  In coordinates, every
distribution \(\partial^2\tau_a/\partial u_i\partial u_j\) vanishes.
Convolution on relatively compact Euclidean balls shows that every local
regularization is affine; compatibility on overlapping balls and passage
to the continuous limit give
\begin{equation}\label{eq:p062-arbitrary-field-affine-trace}
                    \tau_a(u)=c+\sum_{i=1}^n\lambda_i u_i.
\end{equation}

Fix \(i\) and all \(u_j\) with \(j\ne i\).  As \(u_i\to+\infty\), the
Gauss seminorms converge to the Gauss seminorm on the coordinate section
\(T_i=a_i\).  This follows directly from
\[
 \left|\sum_{\nu\in\mathbf N^n}c_\nu(T-a)^\nu\right|_{
      \eta_{a,(\mathrm e^{-u_1},\ldots,\mathrm e^{-u_n})}}
       =\max_\nu |c_\nu|\mathrm e^{-\langle\nu,u\rangle},
\]
first for polynomials and then, by uniform approximation, on a fixed
closed polydisc.  The limiting point lies in \(D^-_{\mathbf R}\), and
continuity of \(h\) gives a finite limit.  Hence \(\lambda_i=0\) in
\eqref{eq:p062-arbitrary-field-affine-trace}.  Repeating this for every
\(i\) shows that \(\tau_a\) is constant.  Letting all coordinates tend
to \(+\infty\) gives its value \(h(a)\).  Because of
\eqref{eq:p062-rho-above-r}, the point corresponding to the originally
chosen \(\mathbf r\) belongs to this trace.  This proves
\eqref{eq:p062-rational-centred-radial-identity} for every allowed
\(\mathbf r\).

If \(a,b\in D^-_{\mathbf R}(K)\), choose
\[
                  |a_i-b_i|\leq r_i<\rho_i<R_i.
\]
Translation invariance of the Gauss norm gives
\(\eta_{a,\mathbf r}=\eta_{b,\mathbf r}\).  The radial identity at the
two centres gives \(h(a)=h(b)\).  Hence the value on every \(K\)-centred
Gauss skeleton is the same, and continuity gives constancy on the closure
of their union.

Now let \(h\) be defined on \(D^+_{\mathbf R}\) and suppose that the
current equation holds on its absolute analytic interior.  For a
\(K\)-point \(a\) and radii \(r_i<\rho_i<R_i\), the set
\(A_{a,\boldsymbol\rho}\) is an ambient open analytic subspace contained
in \(D^+_{\mathbf R}\).  It is therefore contained in the topological
interior of that analytic domain.  Berkovich's relative-interior theorem
and composition formula
\cite[Proposition~1.5.5(i)--(ii)]{Berk93}, together with the
boundarylessness of analytic affine space, place it in the absolute
analytic interior.  The preceding test-form argument applies and proves
the radial identity when every \(r_i<R_i\).  Letting any selected radii
increase to \(R_i\), and using continuity of the Gauss section and of
\(h\), proves it for \(r_i\leq R_i\).  For two \(K\)-points \(a,b\), take
\(0<r_i\leq R_i\) with
\(|a_i-b_i|\leq r_i\), using \(r_i=R_i\) when necessary.  The same common
Gauss point argument finishes the closed case.
\end{proof}

\begin{proof}[Proof of \cref{prop:p062-explicit-base-change}]
Let \(p_L:D^-_{\mathbf R,L}\to D^-_{\mathbf R}\) be the scalar-extension
projection.  Apply \cref{prop:p062-rational-centred-traces} over \(L\) to
\(h_L=h\circ p_L\).  It follows that \(h_L\) is constant on the
\(L\)-rational points.  Rigid points are dense in every strictly affinoid
space \cite[Proposition~2.1.15]{Berk12}; because \(L\) is algebraically
closed, the rigid points of a strict \(L\)-polydisc are precisely its
\(L\)-rational points.  The open polydisc is exhausted pointwise by strict
closed subpolydiscs: the nontrivial divisible group
\(|L^\times|\) is dense in \(\mathbf R_{>0}\), so the subpolydisc radii
may be chosen strictly between the coordinate seminorms of a given point
and the \(R_i\).  Indeed, after applying \(-\log\), this value group is
a nonzero divisible additive subgroup and contains
\(\gamma\mathbf Q\) for some \(\gamma\ne0\).  Therefore the
\(L\)-rational points are dense in
\(D^-_{\mathbf R,L}\), and continuity makes \(h_L\) constant everywhere.
The canonical projection \(p_L\) is surjective
\cite[\S11.1.9]{CLD25}, so \(h\) is constant.
This proves the first item.  The strict closed polydisc is already
affinoid; the closed part of \cref{prop:p062-rational-centred-traces},
the same rigid-point density, and surjectivity of scalar extension prove
the second item.

Assume now that both \(h\) and \(-h\) are CLD-psh in the approximation
sense.  The scalar-extension projection is a morphism of the associated
affined G-tropical spaces: pullback carries every local affine tropical
function to a local affine tropical function
\cite[\S\S12.4.5--12.4.6]{CLD25}.  Functoriality of psh
functions \cite[Proposition~8.1.5]{CLD25} therefore makes \(h_L\) and
\(-h_L\) psh.  The open polydisc is separated, pure \(n\)-dimensional,
and boundaryless, and the two continuous functions are locally
integrable.  Thus their Hessian currents are positive by
\cite[Proposition~8.3.1(a)]{CLD25}.  Since
\[
             \mathrm d'\mathrm d''[-h_L]
                 =-\mathrm d'\mathrm d''[h_L],
\]
both a current and its negative are positive.  Evaluation on positive
test forms shows that both currents vanish.  The first item now applies.
\end{proof}
